\documentclass{article}
\usepackage{amsmath,amssymb,amsthm,mathtools,xcolor}
\usepackage[vocab=unitsquare]{leanify}
% A revision: the changes are printed in colour unless \cleanversion is defined.
\newif\ifmarked
\ifdefined\cleanversion\markedfalse\else\markedtrue\fi
\newcommand{\rev}[1]{\ifmarked{\color{magenta}#1}\else#1\fi}
\newcommand{\revon}{\ifmarked\color{magenta}\fi}
\newcommand{\moved}[1]{\ifmarked\marginpar{\raggedright\scriptsize\color{magenta}#1}\fi}
\newcommand{\CC}{\mathcal{C}}
\newcommand{\de}{\,\mathrm{d}}
\newcommand{\rone}[1]{#1_{1}^{*}}
\newcommand{\rtwo}[1]{#1_{2}^{*}}
\newtheorem{theorem}{Theorem}
\newtheorem{lemma}[theorem]{Lemma}
\begin{document}

Let $\CC$ denote the class of bivariate copulas, that is, distribution functions on $[0,1]^2$ with uniform marginals, and we refer to the literature for background.
The copulas $M(u,v)\coloneqq\min\{u,v\}$, $W(u,v)\coloneqq\max\{u+v-1,0\}$, and $\Pi(u,v)\coloneqq uv$ describe comonotonicity, countermonotonicity, and independence, respectively.
For a copula $C$, Spearman's rho is
\begin{equation}
  \label{eq:rho}
  \rho(C)=12\int_0^1\!\int_0^1 C(u,v)\de u\de v-3,
\end{equation}
as given in \cite[Sec.~5]{nelsen}, a classical measure of concordance.
For a copula $C$, we write $C^\top(u,v)\coloneqq C(v,u)$ for the transpose of $C$.
\rev{For a copula $C$, the survival copula of $C$ is $\widehat C(u,v)\coloneqq u+v-1+C(1-u,1-v)$.
For a copula $C$, the reflections of $C$ in the first and in the second argument are
\[
  \rone{C}(u,v)\coloneqq v-C(1-u,v)
  \qquad\text{and}\qquad
  \rtwo{C}(u,v)\coloneqq u-C(u,1-v).
\]}

\moved{Moved here from the introduction.}%
\rev{The boundary of the region is described by the function} $\Phi:[-1,1]\to[-1,1]$ \rev{given} by
\begin{equation}
  \label{eq:Phi}
  \Phi(r)\coloneqq
  \begin{cases}
    2r+1-\frac34(1+r)^{4/3}, & -1\le r\le0,\\[0.3ex]
    1-\frac34(1-r)^{4/3}, & 0\le r\le1,
  \end{cases}
\end{equation}
\rev{which satisfies} $2r-\Phi(r)=-\Phi(-r)$ for every $r\in[-1,1]$.

\begin{lemma}[\rev{Survival and reflections}]
\label{lem:symmetries}
For every $C\in\CC$,
\begin{equation}
  \rho(\widehat C)=\rho(C),\qquad
  \rho(C^\top)=\rho(C),
\end{equation}
\rev{and the two reflections satisfy}
\begin{equation}
  \revon
  \rho(\rone{C})=-\rho(C),\quad\;
  \rho(\rtwo{C})=-\rho(C),
\end{equation}
\rev{as well as}
\begin{equation}
  \revon
  \rho(\rone{C})+\rho(\rtwo{C})=-2\rho(C).
\end{equation}
\end{lemma}

\end{document}
